\infographicchapter{Dual-Energy X-ray Imaging (DEI)}{dei.png}
\label{ch:infographic-dei}

\section{Overview}

Dual-energy X-ray imaging (DEI)---most often called dual-energy subtraction (DES) chest X-ray---is a specialized digital X-ray method that separates bone from soft tissue in a single exam. Because calcium-rich bone and soft tissues like the lungs absorb low-energy and high-energy X-rays differently, the scanner captures two matched X-ray views (either with two split-second flashes or a special two-layer detector). A computer then combines and subtracts those two views to create three helpful pictures at once: a regular chest X-ray, a ``soft-tissue'' view with the ribs and collarbones digitally removed, and a ``bone-only'' view that highlights bones and calcium deposits.

On a standard chest X-ray, overlapping ribs and collarbones can easily hide small spots in the lungs. DEI helps solve this problem by clearing away bone shadows so doctors can spot subtle lung nodules more easily, tell whether a spot contains harmless calcium (such as an old healed infection) or needs a closer look, and see blood vessel calcium or medical tubes more clearly. Similar dual-energy principles are used in bone density scans (DXA) and dual-energy CT, while chest DEI itself remains a specialized, less commonly used add-on rather than a routine test at every clinic.

The main benefit of DEI is extra clarity over a standard chest X-ray at a tiny fraction of the radiation dose and cost of a chest CT scan. Drawbacks include the need for specialized X-ray equipment, possible blurred double edges if your heart beats or you move between two quick flashes, and slightly grainier subtracted pictures. Safety is much like a standard chest X-ray, with a slightly higher---yet still very low---radiation dose that follows standard pregnancy and dose-safety precautions.

\section{Underlying Physical Principles}
\begin{itemize}
    \item \textbf{Two X-Ray Energy Levels:} DEI captures two matched chest X-rays at different energy strengths---one lower energy (around 60 kVp) and one higher energy (120--140 kVp)---using either two rapid flashes a fraction of a second apart (dual-shot) or a single flash captured by a stacked two-layer detector with a thin copper filter in between (single-shot).
    \item \textbf{How Bone and Soft Tissue React Differently:} Because calcium has heavier atoms (a higher atomic number) than the lighter elements in lung and muscle tissue, bone absorbs low-energy X-rays much more strongly than high-energy X-rays, whereas soft tissues absorb both energies more evenly. Separating the two clears away the ``anatomical noise'' created by overlapping ribs and collarbones.
    \item \textbf{Digital Bone and Tissue Separation:} By comparing how much X-ray energy passed through each spot at the low and high settings, the computer mathematically cancels out either the bones or the soft tissues.
    \item \textbf{Three Views From One Exam:} The software produces three lined-up images---a standard chest X-ray, a bone-removed lung view, and a bone-and-calcium view---while applying gentle smoothing filters to reduce graininess.
\end{itemize}

\section{Imaging Mechanism}
\begin{enumerate}
    \item \textbf{Positioning at the Chest Plate:} You stand upright facing a flat digital X-ray plate in the standard back-to-front (posteroanterior, or PA) position---which keeps your heart close to the detector to avoid artificial magnification---with your shoulders rolled forward so your shoulder blades move out of the way, and the technologist selects dual-energy mode.
    \item \textbf{Split-Second Dual-Energy Capture:} While you hold a deep breath completely still, the machine either fires two rapid X-ray pulses less than a fifth of a second apart (dual-shot) or takes one flash through a two-layer detector (single-shot).
    \item \textbf{Recording Both Energy Views:} The digital detector measures both the low-energy and high-energy X-rays separately and turns them into two crisp digital files.
    \item \textbf{Aligning and Subtracting:} Computer software checks alignment, subtracts out the bone shadows to create the soft-tissue view, subtracts out the soft tissue to create the bone view, and smooths out extra image grain.
    \item \textbf{Reviewing All Three Views Together:} The radiologist views the standard X-ray, the bone-removed lung view, and the calcium-only view side by side on a computer screen.
\end{enumerate}

\section{Key Technical Parameters}
\noindent
\begin{tabularx}{\textwidth}{@{}l X@{}}
\toprule
\textbf{Parameter} & \textbf{Typical Value / Description} \\
\midrule
Physical Principle & Digital comparison and subtraction of low-energy and high-energy X-rays to separate bone/calcium from soft tissue \\
Energy / Signal Source & Digital X-ray tube using two rapid energy flashes (dual-shot) or a two-layer detector with a copper filter (single-shot) \\
Typical Frequency / Energy & About 60 kVp (low energy) and 120--140 kVp (high energy) \\
Spatial Resolution & Very sharp (0.15--0.2 mm pixels), with slightly more grain on the subtracted views \\
Typical Scan Duration & Less than a second for the X-ray capture; 5--10 minutes for the full visit \\
Contrast Agent & None required \\
Ionizing Radiation & Yes; very low dose (about 0.05--0.15 mSv for a chest DEI, slightly higher than one standard chest X-ray) \\
\bottomrule
\end{tabularx}

\section{Common Symptoms That Prompt This Scan}
\begin{itemize}
    \item \textbf{Unclear Spot on a Standard Chest X-Ray:} Often prompted when a routine chest X-ray shows a subtle shadow behind the ribs or collarbones that needs clearer separation of bone and lung tissue.
    \item \textbf{Persistent Cough or Unexplained Chest Symptoms:} Lingering cough, shortness of breath, or mild chest discomfort when doctors want a sharper, low-radiation look at the lungs before considering a CT scan.
    \item \textbf{Sudden Sharp Chest Pain or Shortness of Breath:} Symptoms that raise concern for a small collapsed lung (pneumothorax) near the top of the chest where collarbones can hide air pockets.
    \item \textbf{Localized Rib or Chest-Wall Pain:} Point tenderness over the ribs after a fall or bump when a subtle rib crack is suspected.
\end{itemize}

\section{Clinical Usages}

\subsection{Primary Usages}
\begin{itemize}
    \item \textbf{Finding Hidden Lung Nodules:} Digitally erasing the ribs and collarbones makes small spots near the top of the lungs or behind ribs much easier to see---an important advantage because studies of missed lung cancers on standard chest X-rays found most were hidden behind one or more overlapping bones. Doctors may especially consider it for people at higher lung risk, such as those with a significant smoking history, occupational exposures, or another cancer that can spread to the lungs.
    \item \textbf{Checking for Harmless Calcium:} Showing whether a lung spot lights up on the bone-only view helps doctors tell harmless calcified scars (granulomas) apart from soft-tissue spots that need closer follow-up.
    \item \textbf{Spotting Heart and Chest Calcium:} Highlighting calcium buildup in heart arteries, the aorta, heart valves, the heart sac, or the lung lining on the bone-only view.
    \item \textbf{Checking Airways and the Center of the Chest:} Removing spine and breastbone shadows gives a clearer view of the windpipe, lymph nodes, and areas behind the heart.
\end{itemize}

\subsection{Secondary Usages}
\begin{itemize}
    \item \textbf{Spotting Small Collapsed Lungs (Pneumothorax):} Removing collarbone and upper rib shadows makes it easier to spot a small pocket of trapped air at the top of the lung.
    \item \textbf{Checking Tubes, Stents, and Foreign Objects:} Making catheter tips, stents, pacemaker wires, or swallowed objects stand out clearly on the bone-only view without soft-tissue clutter.
    \item \textbf{Checking Ribs for Cracks or Spots:} Isolating the ribcage helps doctors spot subtle rib fractures or bone spots in the chest wall.
    \item \textbf{Low-Dose Problem Solving:} Serving as a quick, low-radiation, lower-cost helper when a standard chest X-ray is unclear and a full CT scan may not be needed right away.
\end{itemize}

\section{Advantages}
\begin{itemize}
    \item \textbf{Clears Away Rib and Collarbone Shadows:} Removing overlapping bones uncovers subtle lung spots and small air leaks that are easy to miss on a regular chest X-ray; in clinical studies, adding dual-energy views helped radiologists detect more lung nodules (raising sensitivity from 58\% to 67\%) and interpret scans faster.
    \item \textbf{Separates Calcium From Soft Tissue:} Immediately shows whether a lung spot is densely calcified (usually harmless) and highlights calcium in blood vessels or along the lung lining.
    \item \textbf{Much Lower Radiation and Cost Than CT:} Adds extra diagnostic clarity over a standard X-ray at a tiny fraction of the radiation dose, cost, and time of a chest CT scan.
    \item \textbf{Three Useful Views in One Quick Visit:} Produces a normal chest X-ray alongside both bone-only and soft-tissue-only views in a single breath-hold.
\end{itemize}

\section{Disadvantages}
\begin{itemize}
    \item \textbf{Blurring From Movement or Heartbeat:} On two-flash (dual-shot) machines, a heartbeat, slight breath, or shift during the split-second gap can leave double-edge ``ghost'' outlines along the heart, ribs, or diaphragm.
    \item \textbf{Slightly Grainier Pictures:} Subtracting two X-rays naturally adds a bit of graininess to the bone-only and soft-tissue views, especially on single-shot systems or in patients with larger body frames.
    \item \textbf{Slightly Higher Radiation Than a Single X-Ray:} Dual-shot exams deliver roughly twice the radiation of a single standard chest X-ray (around 0.2 mSv versus 0.1 mSv, depending on the scanner)---still a very small dose compared with natural yearly background radiation, though doctors use it only when the extra detail is helpful and take standard care during pregnancy.
    \item \textbf{Still a Flat 2D Picture:} Even with bone removed, DEI cannot provide the full 3D cross-sectional detail of a chest CT scan.
    \item \textbf{Specialized Equipment and Limited Availability:} Requires dedicated dual-energy hardware and software, gives radiologists more images to review, and remains a specialized tool available only at select centers.
\end{itemize}

\section{Contraindications}

\subsection{Absolute}
\begin{itemize}
    \item \textbf{None:} Because DEI is painless, low-dose, and uses no needles or contrast dye, there are no strict medical reasons that forbid it when it is clinically needed.
\end{itemize}

\subsection{Relative}
\begin{itemize}
    \item \textbf{Pregnancy:} Non-urgent dual-energy chest X-rays are usually postponed during pregnancy if a standard single X-ray or radiation-free ultrasound can answer the question.
    \item \textbf{Trouble Holding Your Breath:} On two-flash systems, people who cannot hold their breath still for a few seconds (such as young children or people who are very short of breath) may have blurred subtraction lines, so a single-shot machine or regular chest X-ray is often better.
\end{itemize}

\section{Patient Safety Guidelines}
\begin{itemize}
    \item \textbf{Using It Only When Helpful:} Care teams follow ALARA (as low as reasonably achievable) radiation safety rules, choosing DEI when bone removal or calcium checks truly add value and using child-sized settings for young patients.
    \item \textbf{Tight X-Ray Focus:} Technologists narrow the X-ray beam tightly to the chest and check automatic exposure settings so no extra radiation is given.
    \item \textbf{Coaching a Steady Breath-Hold:} Practicing the breath-hold beforehand helps prevent motion blur on two-flash machines and avoids needing to repeat the X-ray.
    \item \textbf{Pregnancy Screening:} Staff check for possible pregnancy before the scan. Although a chest X-ray sends very little scatter toward the pelvis, teams confirm clinical need and keep the beam tightly focused.
    \item \textbf{Regular Equipment Checks:} Clinics routinely calibrate the detectors and check radiation levels so tissue separation stays accurate without excess dose.
\end{itemize}

\section{Prerequisite Tests}
\begin{itemize}
    \item \textbf{No Blood Tests Needed:} DEI is completely noninvasive and uses no contrast dye, so no blood work, kidney checks, or fasting are needed.
    \item \textbf{Pregnancy Check:} Patients who could be pregnant are asked about pregnancy status before any X-ray exam.
    \item \textbf{Review of Past Chest Scans:} Doctors check prior chest X-rays or CT scans to make sure dual-energy imaging will add helpful information.
    \item \textbf{Breath-Hold and Metal Check:} Staff confirm you can hold a brief breath steadily and make sure all necklaces, snaps, and metal objects are cleared from the chest area.
\end{itemize}

\section{Preparing for the Scan}
\begin{itemize}
    \item \textbf{Remove Chest Metal and Accessories:} Take off necklaces, piercings, underwire bras, metal zippers, buttons, and any removable sticky ECG patches from your chest, because metal objects can look like false calcium spots on the bone view.
    \item \textbf{Wear a Plain Gown:} Change into a metal-free hospital gown or wear a thin, plain shirt without thick seams or metallic prints.
    \item \textbf{Eat and Take Medicines Normally:} No fasting is required, and you can take all your usual medications before the exam.
    \item \textbf{Practice Holding Your Breath:} Be ready to take a full breath in and hold completely still for 2--4 seconds so the low- and high-energy pictures line up cleanly without ghosting.
    \item \textbf{Share Pregnancy or Implant Details:} Let the technologist know if you might be pregnant and mention any implanted chest hardware (like a pacemaker, stent, or surgical clips) so the radiologist recognizes them right away.
\end{itemize}

\section{During Scan}
\begin{itemize}
    \item \textbf{Standing Still at the Plate:} You will stand upright facing the flat detector plate with your chin slightly raised, shoulders rolled forward, and hands on your hips (or sit upright if standing is hard). Staying completely still keeps both energy views aligned.
    \item \textbf{Simple Breath-Hold:} The technologist will ask you to take a deep breath in and hold it steadily for about 2--4 seconds until they tell you to breathe normally.
    \item \textbf{What You Will Feel and Hear:} The test is completely painless and feels just like a regular chest X-ray. You will feel the cool plate against your chest and hear one or two quick beeps or clicks.
    \item \textbf{Staying in Touch With Staff:} The technologist watches you through a window and speaks over an intercom. Tell them right away if dizziness, chest pain, or shortness of breath makes it hard to stand or hold your breath.
    \item \textbf{Quick Visit:} The X-ray itself takes only a second or two, and the whole appointment---including getting positioned and checking that the images are crisp---usually takes 5--10 minutes.
\end{itemize}

\section{Normal Appearance}
\begin{itemize}
    \item \textbf{Standard Chest View:} Looks just like a normal chest X-ray, with clear, air-filled lungs, a normal heart outline, and intact ribs.
    \item \textbf{Bone-Removed Soft-Tissue View:} Ribs, collarbones, shoulder blades, and spine fade away cleanly, leaving clear views of the lungs, lung vessels, windpipe, and top of the lungs with no unusual spots.
    \item \textbf{Bone-and-Calcium View:} Shows the ribs, collarbones, breastbone, and spine crisply against a dark background, with intact bone edges and no abnormal calcium spots in the lungs or blood vessels.
\end{itemize}

\section{Observing for Abnormal Patterns}
\begin{itemize}
    \item \textbf{Newly Visible Soft-Tissue Lung Spots:} A small spot that appears clearly once the ribs are removed---and vanishes on the bone-only view---is a non-calcified lung nodule, which usually leads to a chest CT scan to check it closely.
    \item \textbf{Harmless Calcified Spots:} A lung spot or lining patch that stays bright on the bone-only view and disappears on the soft-tissue view contains dense calcium, which strongly points to a harmless healed scar (granuloma) or calcified plaque.
    \item \textbf{Heart and Vessel Calcium:} Bright curved or ring-shaped lines on the bone-only view clearly show calcium in the heart arteries, aorta, valves, or sac around the heart.
    \item \textbf{Small Air Leaks and Rib Injuries:} Removing collarbone shadows makes small pockets of trapped air at the top of the lung easier to spot, while the bone-only view highlights subtle rib cracks, bone spots, or medical tube tips.
    \item \textbf{Motion Ghosting and Graininess:} On two-flash machines, a heartbeat or slight movement can leave bright-and-dark double edges around the heart, ribs, or diaphragm that look like trapped air or thickened lining. Extra graininess in larger patients can also mimic faint lung haze.
\end{itemize}

\section{Dealing with Uncertainties}
\begin{itemize}
    \item \textbf{Double-Edge Motion Blur:} Movement or heartbeats between the two quick flashes can create ``ghost'' edges that mimic a small collapsed lung or chest lining change.
    \item \textbf{Extra Image Grain:} Subtracted pictures naturally have more grain than standard X-rays---especially on single-shot machines or in larger patients---which can hide faint spots or mimic subtle lung haze.
    \item \textbf{Partially Calcified or Overlapping Spots:} Some tumors contain partial calcium, and normal skin moles, nipple shadows, or rib cartilage calcium can look unclear on a flat 2D image.
    \item \textbf{Checking All Three Views Together:} Radiologists always compare the standard X-ray, soft-tissue view, and bone view together to spot ghosting or skin shadows. A scan is repeated only if heavy motion makes it unreadable.
    \item \textbf{Comparing With Old Scans or Doing a CT:} Comparing with past chest X-rays helps confirm that a harmless-looking spot has not changed. Any new non-calcified lung nodule or unclear chest finding is checked further with a chest CT scan.
\end{itemize}

\section{Augmenting Methods}
\begin{itemize}
    \item \textbf{Chest CT Scan:} The definitive 3D follow-up test for any newly spotted non-calcified lung nodule, unclear center-chest shadow, or small collapsed lung found on DEI.
    \item \textbf{Dual-Energy CT (DECT):} Uses the same dual-energy idea in full 3D to map lung blood flow for clots and analyze calcium versus contrast dye in chest spots.
    \item \textbf{Digital Chest Tomosynthesis:} Takes low-dose layered chest views to pinpoint small lung spots or rib cracks without jumping straight to a full CT scan.
    \item \textbf{AI Bone-Removal and Detection Software:} Computer programs help smooth out graininess, correct tiny motion shifts, and flag subtle lung spots for the radiologist to check.
    \item \textbf{PET/CT and Biopsy:} If a CT-confirmed lung nodule grows or looks suspicious, a PET/CT scan or tissue biopsy provides the final answer.
\end{itemize}

\section{Comparison with Alternative Modalities}
\begin{itemize}
    \item \textbf{Compared With a Standard Chest X-Ray:} DEI finds hidden lung spots behind ribs more reliably and immediately tells calcified scars from soft-tissue spots, with only a tiny increase in radiation.
    \item \textbf{Compared With a Low-Dose Chest CT:} Chest CT gives full 3D views and remains the gold standard for lung cancer screening and nodule follow-up, while DEI is faster, cheaper, and uses much less radiation as a first-step helper.
    \item \textbf{Compared With Chest Tomosynthesis:} Tomosynthesis gives layered depth slices for locating small spots, whereas DEI uniquely separates calcium from soft tissue in a single quick breath-hold.
\end{itemize}

\section{Follow-up Steps}
\begin{itemize}
    \item \textbf{Go Right Back to Your Day:} No recovery time or special care is needed after DEI. Once the technologist confirms the pictures are sharp and free of motion blur, you can leave right away.
    \item \textbf{Side-by-Side Expert Review:} A radiologist reviews all three views alongside any older chest scans to determine whether a spot is a true lung nodule, a harmless calcified scar, or a motion artifact.
    \item \textbf{Follow-Up CT When Needed:} Newly discovered non-calcified lung spots or unclear chest shadows are typically followed up with a dedicated chest CT scan for full 3D clarity.
    \item \textbf{Tailored Monitoring:} Clearly harmless calcified scars need no further testing, while CT-confirmed nodules are monitored on a standard schedule (such as Fleischner Society or Lung-RADS guidelines).
\end{itemize}

\section{Key References}
\begin{itemize}
    \item Kuhlman JE, Collins J, Brooks GN, Yandow DR, Broderick LS. Dual-Energy Subtraction Chest Radiography: What to Look for Beyond Calcified Nodules. \textit{RadioGraphics}. 2006;26(1):79--92.
    \item Gilkeson RC, Sachs PB. Dual Energy Subtraction Digital Radiography: Technical Considerations, Clinical Applications, and Imaging Pitfalls. \textit{Journal of Thoracic Imaging}. 2006;21(4):303--313.
    \item MacMahon H, Naidich DP, Goo JM, et al. Guidelines for Management of Incidental Pulmonary Nodules Detected on CT Images: From the Fleischner Society 2017. \textit{Radiology}. 2017;284(1):228--243.
    \item American College of Radiology (ACR). \textit{ACR--SPR--STR Practice Parameter for the Performance of Chest Radiography}. Reston, VA: ACR; 2022.
    \item Manji F, Wang J, Norman G, Wang Z, Koff D. \textit{Comparison of dual energy subtraction chest radiography and traditional chest X-rays in the detection of pulmonary nodules}. National Library of Medicine (PubMed Central). Available at: \url{https://pmc.ncbi.nlm.nih.gov/articles/PMC4775247/}.
    \item Radiopaedia.org. \textit{Chest radiograph}. Radiopaedia. Available at: \url{https://radiopaedia.org/articles/chest-radiograph}.
\end{itemize}

\section{Further Reading}
\begin{itemize}
  \item Wikipedia contributors. \textit{Radiography}. Wikipedia, The Free Encyclopedia. Available at: \url{https://en.wikipedia.org/wiki/Radiography}.
\end{itemize}
